% =====================================================================
%  HIKET — manuscript STORY-STRUCTURE outline (scaffold only)
%  Placeholders + brief notes. Headings to be collapsed/merged later.
%
%  PUBLIC narrative schema = Schimel OCAR (Opening / Challenge / Action /
%  Resolution), used as the four \part divisions below.
%
%  PRIVATE scaffold (NEVER surfaced in the text — author reference only):
%    Shadow      = the observed accumulation->saturation is precisely what a
%                  steady-state start CANNOT produce (an equilibrium start is
%                  already "saturated"); this triggered transient init
%    Persona     = good R-hat / physical-looking bias masking artefacts
%    Mentor      = published priors/lineage: real but boreal-blind, then freed
%    Trickster   = sigma_input, the shape-shifter that reveals the real
%                  treasure (input/history uncertainty)
%    Unconscious = the unrecoverable pre-observation history; the humus pool
%                  that "remembers" equilibrium (deep time made physical)
%    Self        = the fair, equal-DoF, exactly-integrated, homogenized frame
%    Boon        = a method + a relocated question, not a verdict on Finland
%
%  Protagonist (held constant): the INITIALIZATION PROBLEM — how to start a
%  non-equilibrium simulation with no history — with transient init as the move.
%
%  RECURRING THREAD — Bayesian inference (tagged % [Bayes] at each touchpoint):
%  state of the art, routinely applied, only mildly novel here — but it is the
%  machinery that MAKES the solution possible: it turns the unknown initial
%  state from a fixed assumption into a CALIBRATED parameter (sigma_init) whose
%  uncertainty propagates into both the initial state and the forward run.
%  Surfaces 4x: Challenge (fixed-value trap vs. calibrate-it), Action (the
%  enabling machinery + anchor propagation), Resolution (the lens that makes
%  non-identifiability/sigma_input visible), Outlook (priors from reconstructed
%  history). Keep it a thread, not a section — it must not crowd the protagonist.
%
%  RECURRING THREAD — historical perspective (tagged % [History]): a quick
%  history of (Finnish) forest management is what PHYSICALLY justifies the
%  non-equilibrium — the violation of the equilibrium assumption. General mention
%  in the Intro (own small subsubsection) -> detailed examination in the
%  Discussion -> closes the Conclusions (refinement needs historical data).
% =====================================================================

\documentclass[11pt]{article}

\usepackage[margin=2.5cm]{geometry}
\usepackage{amsmath}
\usepackage{titlesec}
% ---------------------------------------------------------------------
%  MODULAR APPENDICES.
%  Each appendix lives in appendices/*.tex as a `subfiles` document. It can be:
%    (a) compiled STANDALONE  -> `pdflatex appendices/appendix_xxx.tex`
%        (borrows this file's preamble automatically), or
%    (b) compiled TOGETHER    -> pulled in here via \subfile{...} (see the
%        Appendices part near the end). Comment a \subfile line to drop a module.
%  This keeps each methods insert self-contained and separately shareable.
% ---------------------------------------------------------------------
\usepackage{subfiles}
% ---------------------------------------------------------------------
%  OPTIONAL STRATIFIED-CALIBRATION STRAND (rendered in red).
%  A self-contained future-work line that can be ADDED if data/time allow, or
%  DROPPED entirely with a single switch while the whole storyline stays in place
%  (just smaller). Set \stratfalse to remove every \ifstrat...\fi block below;
%  while \strattrue they render in red so the strand is visually separable.
% ---------------------------------------------------------------------
\usepackage{xcolor}
\newif\ifstrat
\strattrue        % <-- \stratfalse drops the entire stratified-calibration strand
\newcommand{\stratnote}[1]{\ifstrat{\color{red}#1}\fi}  % inline red, droppable

% Number and list every heading level down to subparagraph.
\setcounter{secnumdepth}{5}
\setcounter{tocdepth}{5}

% Make \paragraph and \subparagraph display as block headings (outline view).
\titleformat{\paragraph}[block]
  {\normalfont\normalsize\bfseries}{\theparagraph}{1em}{}
\titleformat{\subparagraph}[block]
  {\normalfont\normalsize\bfseries\itshape}{\thesubparagraph}{1em}{}
\titlespacing*{\paragraph}{0pt}{\baselineskip}{0.3\baselineskip}
\titlespacing*{\subparagraph}{0pt}{\baselineskip}{0.3\baselineskip}

\title{[PLACEHOLDER TITLE]\\
  \large Starting from the past: transient initialization and the
  input-history problem in soil carbon models}
\author{Lorenzo Menichetti \and [co-authors TBD]}
\date{Draft outline --- \today}

\begin{document}
\maketitle

\begin{abstract}
\noindent[PLACEHOLDER ABSTRACT.] One paragraph: steady-state initialization
is a near-universal but rarely examined convention; an equilibrium start is
already at its saturated end-state, so it cannot reproduce an accumulating
sink that is only now beginning to saturate --- the very dynamic observed in
Finnish forest soils. We test a transient alternative across six models under
one fair frame; it captures the accumulation-then-saturation as a proof of
concept, but relocates the binding uncertainty onto litter inputs ---
especially their unknown history --- pointing to historical, interdisciplinary
reconstruction as the real frontier.
\end{abstract}

\tableofcontents

% =====================================================================
\part{Opening}
% Schimel O: broad context + stakes. Lead with the GENERAL convention;
% Finland arrives as the vivid instance. (Generalization-first, confirmed.)
% =====================================================================

\section{Introduction}

\subsection{The convention: steady state at $t_0$}
A near-universal default in SOC modelling. State it plainly as common
practice, not yet as a problem.

\subsubsection{What the assumption actually claims}
That the system's past equals its present equilibrium (inputs $\approx$
outputs at the start). Make the hidden content of the convention explicit.

\paragraph{Why it is attractive}
Convenient, analytically clean, needs no historical forcing. The reasons it
became default.

\paragraph{Why it is rarely examined}
It is invisible precisely because it is standard. Set up the reader to
recognize their own practice.

\subsubsection{The non-equilibrium reality}
Real ecosystems are transient; they carry the imprint of disturbance,
land-use and productivity history.

\paragraph{An equilibrium start is already at the end}
A steady-state start sits at zero net change --- by construction already at
its saturated end-state. It forecloses the accumulation phase before the
simulation begins. Foreshadow the capture-failure without resolving it.

\subsubsection{Managed landscapes are rarely at equilibrium}
% [History] — general statement here; Finnish forest history is the worked
% example, examined in detail in the Discussion.
Land-use and management histories keep soils in transient states; equilibrium
is the exception, not the rule. Flag Finnish forest history as the concrete
case that justifies the violated assumption --- stated generally here, examined
in detail in the Discussion.

\subsection{Why this matters}

\subsubsection{Scientific stakes}
Initialization choices propagate into every downstream projection and into
model intercomparison.

\subsubsection{Policy stakes: the motivating case}
Finland reports its forest soil as a sink; climate policy leans on it.
Introduce Finland as the instance, not the subject.

\paragraph{The observed dynamic}
Accumulation (VMI8$\to$Biosoil), then recent hints of saturation (Komeetta,
under review). A sink caught in the act of saturating. State the phenomenon
the model must reproduce.

\subsection{Scope and contribution (preview)}
One short paragraph previewing the move: a transient initialization, tested
fairly across a model ladder. Keep it a promise, not the argument.
% CONTRIBUTION DELTA --- state explicitly (the ``relatively new, applied'' register;
% reviewers respect it if precise, and ``spin-up is old'' is the reflex to pre-empt).
% NOT transient init per se, but:
%   (i)  the unknown historical start made a CALIBRATED, uncertainty-PROPAGATING
%        parameter (sigma_init) when NO historical forcing data exist --- not a fixed
%        spin-up nor a prescribed land-use history;
%   (ii) applied HOMOGENEOUSLY across a structural ensemble as a fair-frame device
%        (lets divergence be attributed to structure vs. initialisation artefact);
%   (iii) at national-NFI scale (~350-400 plots, per-plot).
% The insight it buys: inputs/forcing history bind; structure does not.

% =====================================================================
\part{Challenge}
% Schimel C: the knowledge gap and the question. The Shadow lives here.
% =====================================================================

\section{The gap and the question}

\subsection{The saturation an equilibrium start cannot show}
The triggering insight: to reproduce a sink that is saturating, the model must
first reproduce the accumulation it is the tail of --- and an equilibrium
start, already saturated, cannot. This is what motivated transient init.

\subsubsection{Why this generalizes beyond Finland}
Any non-equilibrium system started at steady state forecloses its own
transient. Lift the problem off the case study.

\subsubsection{Artefact or structure?}
Do divergent saturation projections reflect genuine model structure, or the
initialization and calibration convention? The question the fair frame will
answer.

\subsection{The methodological gap}
No established, calibratable way to initialize a transient run when
historical forcing data are absent.

\subsubsection{The fixed-value trap}
% [Bayes]
Treating the unknown initial state as a single fixed number (equilibrium, or a
guessed value) hides a strong, untested assumption inside the result.

\subsubsection{Calibrate it, do not assume it}
% [Bayes] — first appearance of the thread; the gap is what Bayes will close
The alternative: make the historical start a parameter with a prior and let
the data and its uncertainty speak. Plant the thread; pay it off in Methods.

\subsection{Objectives and hypotheses}
\begin{itemize}
  \item[H1] Transient initialization lets the models reproduce the
        accumulation-then-saturation a steady-state start cannot. [placeholder]
  \item[H2] Residual model differences under a fair frame are structural and
        modest. [placeholder]
  \item[H3] The binding uncertainty migrates to inputs, not kinetics.
        [placeholder]
\end{itemize}

% =====================================================================
\part{Action}
% Schimel A: methods + results. The Self (fair, whole frame) is built here.
% =====================================================================

\section{Materials and methods}
% Detailed derivations are offloaded to the modular appendices (see the
% Appendices part); the M&M body states each choice and points to its appendix.

\subsection{Study system and data}

\subsubsection{Plots and SOC campaigns}
Finnish NFI plots; three temporally separated SOC campaigns as targets.
Brief.

\subsubsection{Litter inputs}
Plot-level AWEN-fractionated inputs; note the population-mean rise then
plateau (load-bearing later).

\subsubsection{Climate forcing}
Gridded daily climate; per-model temporal aggregation.

\subsection{The model ladder}
Six models on two crossed axes (pool complexity; climate integration), so no
conclusion rests on one structure.
% MODEL ROLES (state up front): Yasso is the PRIMARY / operational model --- it is
% the Finnish GHG-inventory model, so it is the one that matters for the application.
% SP1/TP2/TP3 are a structural ROBUSTNESS ladder: hypothesis-testing that the
% transient-init + input-uncertainty conclusions survive across model structure,
% NOT a beauty contest for the ``best'' model. This reframes the complexity axis as
% robustness/sensitivity rather than an optimal-complexity search.

\subsubsection{Pool-complexity axis}
SP1 $\to$ TP2 $\to$ TP3 $\to$ Yasso07. One sentence each. [placeholder]

\subsubsection{Climate-integration axis}
Yasso07 $\to$ Yasso15 $\to$ Yasso20. One sentence each. [placeholder]

\subsection{Bayesian calibration: the enabling machinery}
% [Bayes] — the thread's payoff. State of the art, routinely applied; here it
% is what makes transient init possible. Implementation detail, but fundamental.
The inference framework that lets the unknown start be calibrated rather than
assumed. Keep it brief and matter-of-fact; its role is enabling, not headline.

\subsubsection{Why a Bayesian frame is the natural fit}
Priors encode what little is known about the past; the posterior carries
uncertainty forward instead of collapsing it to a point.

\subsubsection{Uncertainty that propagates, not a likelihood patch}
$\sigma_{\mathrm{init}}$ is a prior on the historical state; its spread flows
into both the initial condition and the forward dynamics. Distinguish from a
mere error term added at the end.

\subsection{Transient initialization (the protagonist)}

\subsubsection{Pushing the equilibrium assumption back in time}
Initialize at analytical steady state in the past, run transiently forward to
$t_0$ so the imprint dissipates and the model arrives at $t_0$ still
accumulating.

\paragraph{The fast/slow relaxation argument}
Fast pools fully relax; only the slow humus pool retains memory --- it
``remembers'' the equilibrium start. Why the pre-run works, and where its
memory lives.

\subsubsection{Calibrating the historical anchor}
% [Bayes]
$\sigma_{\mathrm{init}}$ as a Bayesian-calibrated strength on the historical
start, not a fixed choice --- the concrete place the thread meets the
protagonist. Full pre-run construction in Appendix~\ref{app:transient}.

\subsection{A fair frame for intercomparison}
The design principle: any residual difference must be attributable to
structure alone.

\subsubsection{Shared data, likelihood, forcing, filtering}
Brief enumeration.

\subsubsection{Exact integration in every model}
All six integrate their linear system exactly; differences are structural,
not numerical. Solvers and the discarded Euler path in
Appendix~\ref{app:exact}.

\subsubsection{Homogenized priors: method, not number}
Tiered scheme; identical in construction/scale/DoF, not in raw values. Full
parameter-class $\times$ prior-criteria tables in Appendix~\ref{app:priors}.

\subsubsection{Degrees of freedom, set by external information}
Flow fractions are equal-DoF (all free) in every multi-pool model --- the
fair-frame core. Rates differ BY DESIGN, and the reason is information, not
inconsistency:
% Yasso FIXES its base rates (published, authoritative -> preserves each pool's
% chemical definition) and spends its freedom on fractions + climate + size.
% SP1/TP2/TP3 CALIBRATE their rates because no authoritative rate constants exist to
% fix them to (they have priors, but no external value worth pinning). So the ladder
% spans externally-anchored (Yasso) -> fully-calibrated (simple): identifiability
% tracks how much external information a structure carries. State openly --- a
% careful reader WILL notice the rate asymmetry; framed this way it is a deliberate,
% information-driven choice, and it doubles as an illustration of the equifinality
% theme (see Discussion excursus).
Pin what external knowledge lets you pin; calibrate the rest.
\stratnote{[strand] This convention is also what makes the stratified extension
Yasso-only: with rates fixed, Yasso's stratifiable knobs are inputs + the
humus-formation fraction --- see the optional strand in the Outlook.}

\subsection{Calibration and inference}

\subsubsection{Error model}
Multiplicative-normal likelihood. Flag the asymmetry briefly (paid back in
Discussion).

\subsubsection{Physically-bounded input priors}
% Now a METHODS choice (implemented for the production re-calibration), not only
% an Outlook refinement. The effective litter flux is bounded to the boreal
% NPP envelope, homogeneous across models, via a bounded re-parameterisation.
The effective litter flux is bounded to a boreal NPP envelope, homogeneous across
models, so the input multiplier cannot manufacture unphysical carbon; the
auxiliary $\sigma_{\mathrm{init}}$/$\sigma_{\mathrm{input}}$ pair is
re-parameterised rather than truncated. Rationale, envelope and mechanism in
Appendix~\ref{app:input-bounds}. [placeholder]

\subsubsection{Sampler and diagnostics}
MCMC settings; convergence and identifiability diagnostics used.

\subsubsection{Holdout split and validation}
Plot-level holdout excluded from the likelihood; separate predictive metrics.
[placeholder]

\subsubsection{Software, compute and reproducibility}
Engine, samplers, exact-integration back-ends; fixed seeds; HPC production run.
[placeholder]

\subsection{Predictive and residual analysis}
Posterior-predictive trajectories; residual analysis against site covariates.

\section{Results}

\subsection{Capturing accumulation-then-saturation}
The headline: transiently-initialized models reproduce the rise and its
incipient bend-over; the equilibrium-started flat/declining shape was the
start, not the structure.

\subsubsection{Trajectory shape vs.\ level}
Shape change, not just offset. [placeholder figure]

\subsubsection{Match to campaign means}
2006 and 2024 means tracked. [placeholder]

\subsection{Structural differences under the fair frame}
Real but modest skill spread across the six. [placeholder table]

\subsubsection{Where simple beats complex}
TP2/TP3 track shape best; stiffer Yasso family worst. Brief.

\subsection{What the data can and cannot constrain}
Set up the Discussion: kinetics weakly constrained; inputs do the work.

\subsubsection{Convergence vs.\ identifiability}
Good R-hat coexisting with non-identified fractions. State the apparent
paradox, resolve in Discussion.

\subsubsection{Input multiplier behaviour}
$\sigma_{\mathrm{input}}$ does different work in different models. The hinge
result.

\subsection{Residual structure}
RF importance; stand basal area as the leading, diffuse, shared driver.
[placeholder figure]

% =====================================================================
\part{Resolution}
% Schimel R: discussion + significance + the opened frontier. The Boon.
% =====================================================================

\section{Discussion}

\subsection{Capturing saturation requires remembering the accumulation}
Resolve the Challenge: the difficulty was the equilibrium start, not the soil;
a transient start lets every model represent the observed dynamic. The central
payoff.

\subsubsection{Proof of concept, not a verdict on Finland}
What the result does and does not license. Calibrate the claim.

\subsubsection{Generalization to the model class}
Wherever steady-state init meets a non-equilibrium system, it forecloses the
transient the data demand. Lift off Finland.

\subsection{The binding uncertainty is the inputs (the $\sigma_{\mathrm{input}}$ reveal)}
% Distinct beat (chosen: clearer over woven). Trickster pays off here.
The data cannot resolve kinetics; the input multiplier absorbs what structure
cannot. So the frontier is inputs, not soil chemistry.
% CENTRALITY: not a side story --- it is the ANSWER to the complexity question
% (structure is a weak lever BECAUSE inputs dominate). Give it Discussion-central
% status, not a sidebar.
% CLOSED LOOP: diagnosed on the unbounded run -> bounded in Methods -> consequence
% on the bounded re-run (whether TP2/TP3 take a visible R^2 hit, or the misfit
% relocates onto the now-pinned kinetics). The bound is what gives the reveal teeth.
% HEDGE (honesty): say ``consistent with input/forcing error'', NOT ``proven'' ---
% the current global design cannot exclude local kinetic variation (see excursus +
% the optional stratified strand, which is exactly the fair test of that).

\subsubsection{The multiplier crosses a physical line (quantitative payoff)}
% The Trickster's tell: sigma_input is not a nuisance, it has a physical referent.
% Effective flux = sigma_input * raw litter (raw ~2.5 tC/ha/yr, shared by all six).
% Spread: SP1 0.8, Yasso07 3.5, Yasso15 4.8, Yasso20 7.5 (all physical, <= ~9 NPP
% ceiling) BUT TP2 34, TP3 50 tC/ha/yr -- 13-20x, above boreal NPP -> IMPOSSIBLE.
% The good R^2 of TP2/TP3 is bought by manufacturing carbon. Same finding as the
% TP3 oscillation (fast pool force-fed). Reframe: sigma_input distance from ~1 is a
% STRUCTURAL-ADEQUACY GAUGE, interpretable only once the prior encodes the physical
% litter bound. Homogenizing sigma_input as a weak Tier-3 nuisance (log-SD 0.5)
% threw away external litter information -- the one prior class that HAS a referent.

\subsubsection{Historical inputs: the unrecoverable past}
No forcing data pre-observation; $\sigma_{\mathrm{init}}$ stands in for an
entire history, held in the slow pool that remembers it. The deepest
uncertainty.

\subsubsection{Belowground inputs and other ramifications}
Root/rhizodeposition inputs poorly known; a further axis of the same problem.

\subsection{A short excursus: why this class of models is non-identifiable}
% ~1/3 page. Side story that connects to the sigma_input reveal.

\subsubsection{The structural argument}
First-order kinetics (and richer topologies) let parameters interact so the
data constrain net fluxes out of a pool, not the individual splits.

\paragraph{Sketch}
For a linear pool system the observable stocks depend on the parameters only
through lumped combinations; distinct parameter sets give identical
predictions (parameter-space ridges). Two equations of placeholder math.
Full argument in Appendix~\ref{app:nonident}. [placeholder]

\paragraph{Signature in practice}
% [Bayes] — the lens: the Bayesian frame makes non-identifiability visible and
% honest (posterior vs prior, KL) rather than hiding it in a point estimate.
Prior-pinning, near-perfect anti-correlation ridges between same-pool
fractions, low information gain (KL of posterior vs prior). The Bayesian frame
is what lets us see this rather than hide it in a fitted point. Tie to Results.

\paragraph{Why it is not a defect to be ``fixed''}
It is intrinsic to the model class and the data content (few observations per
plot); it bounds what calibration can ever claim.

\paragraph{Managing it by convention, transparently}
We do not claim to resolve the rate$\leftrightarrow$fraction equifinality --- we
MANAGE it: where authoritative rate values exist (Yasso) we pin the chemical axis
and calibrate the partitioning; where they do not (simple models) we calibrate and
let the equifinality show. An explicit convention, not an identification claim ---
and the candour is what keeps this section credible rather than defensive. Ties
back to ``degrees of freedom set by external information'' in Methods.

\subsection{The over-prediction is a property of the likelihood}
% Persona/candor beat.
Slight, systematic over-prediction follows from the multiplicative-normal
error model, not from biogeochemistry. Only the bias is likelihood-sensitive;
rankings are not.

\subsection{Limitations}

\subsubsection{The coarse pre-observation phase}
Linear litter interpolation under constant climate; high, deliberate
uncertainty.

\subsubsection{The flat-vs-rising litter tension}
Plot-level inputs vs.\ national reconstructions; present data cannot settle
it. State openly.

\subsection{A managed landscape out of equilibrium: Finnish forest history}
% [History] — the detailed examination promised in the Intro; justifies the
% non-equilibrium and the rising-input trajectory; feeds the reconstruction call.
The concrete history behind the violated equilibrium assumption: why Finnish
forest soils were accumulating, not at steady state, across the observation
period. Pays off the Intro's general statement.

\subsubsection{From exploitation to managed growth}
A sketch of the management eras --- heavy historical use, then regeneration and
intensified management, rising growing stock through the 20th century.
[placeholder --- sources to add]

\subsubsection{The imprint on litter inputs and stocks}
Rising biomass and shifting management make the input history non-stationary
--- precisely the signal an equilibrium start erases. Connect to the
$\sigma_{\mathrm{input}}$ reveal and the flat-vs-rising litter tension.

\subsubsection{Why this is the rule, not a Finnish peculiarity}
Most managed and recovering landscapes carry comparable histories; the lesson
generalizes. Tie back to the Intro's general statement.

\subsection{Outlook: toward a historical perspective on inputs}
The forward call. The real next step is reconstructing historical (and
belowground) input regimes.
% TWO-PRONGED FRONTIER, one lesson (add structure to the FORCING, not the kinetics):
%   (temporal) reconstruct the input HISTORY --- replace the linear 1917->1985 guess;
%   (spatial/class) stratify the input LEVEL by site class --- the red optional strand.
% The physical envelope UNIFIES both: it bounds sigma_init (the 1917/history input)
% AND sigma_input (the contemporary input) under one constraint, so ``reconstruct the
% input history'' literally means doing better than our linear guess at BOTH ends of
% the pre-run. Note: Komeetta 2024 = a 3rd SOC obs/plot -> won't rescue kinetic
% identifiability (12 fractions vs 3 points), but strengthens the input/init signal
% that the stratified strand leans on.

\subsubsection{From diagnosis to method (now folded into M\&M)}
% The physical input bound was the ``refinement'' beat --- it is now IMPLEMENTED as
% a Methods choice (M&M ``Physically-bounded input priors'' + App.~\ref{app:input-bounds}),
% so it is no longer future work. What remains genuinely FUTURE is the input HISTORY
% (temporal reconstruction) and belowground inputs.
The bound itself has moved to Methods; the remaining frontier is the input
\emph{history} (temporal) and belowground inputs. See Appendix~\ref{app:input-bounds}.

\subsubsection{An interdisciplinary, source-diverse effort}
Reconstruction needs evidence beyond the usual datasets --- grey literature
and non-traditional historical sources. Frame as a research program.

\subsubsection{Initialization as a calibratable, evidence-based step}
% [Bayes] — the thread closes: reconstructed history becomes the prior.
Replace the equilibrium default with priors grounded in reconstructed history,
calibrated under uncertainty. The thread closes here: better history $\to$
sharper prior $\to$ a start that is assumed less and inferred more. The
generalizable recommendation.

% =====================================================================
%  OPTIONAL STRATIFIED-CALIBRATION STRAND  (red; drop entirely with \stratfalse)
%  Self-contained: removing it leaves the storyline whole, just smaller.
% =====================================================================
\ifstrat
{\color{red}
\subsection{[Optional strand] Locally stratified calibration (Yasso only)}
A droppable extension --- include only if data + time allow; the storyline stands
without it. It is the SPATIAL counterpart of the historical-reconstruction prong:
both add structure to the \emph{forcing}, never to the kinetics.

\subsubsection{Scope --- the operational model only}
Stratify Yasso (the operational model); SP1/TP2/TP3 stay global robustness anchors.
Consistent with the rate convention: Yasso already fixes its published rates, so its
free knobs (inputs + fractions) are the ones available to stratify. Triggered only
``if complexity emerges as needed'', not by default.

\subsubsection{What to stratify --- inputs and the humus-formation fraction}
The physically-bounded input $\sigma$ and the humus-formation fraction (the
stabilisation control) --- NOT the rates (fixed) nor the inter-pool splits
(non-identified). Both carry external, class-varying referents; humus formation is
a near-net stabilisation flux the total stock genuinely constrains, so of the
fractions it is the one worth stratifying.

\subsubsection{RF-guided selection, holdout-evaluated}
Use the residual random forest to CHOOSE the stratifying classes (its OOB $R^2$ is
the ceiling of covariate-recoverable local structure); then measure the stratified
model's gain on the HELD-OUT plots --- non-negotiable, or selecting and scoring on
the same data makes the improvement circular. Report ``RF says up to $X\%$ is
recoverable; the stratified mechanistic model recovers $Y\%$ on holdout.''

\subsubsection{Hierarchical, predictive-robustness framing}
Partial pooling: data-poor strata shrink to the global value (no minimum-$N$
problem); the between-stratum variance is itself a result. Headline = improved,
more ROBUST posterior predictions even where stratum parameters stay weakly
identified --- identifiability and predictive skill are distinct, and robustness is
what the GHG-inventory use needs.

\subsubsection{Equifinality as a measured quantity}
Bounded inputs vs.\ free humus formation compete for the local variance; the
physical bound PARTIALLY breaks their equifinality (one competitor has walls, the
other does not), so the learned variances softly attribute local variance to
input vs.\ stabilisation origin --- cross-checked by the RF driver identity
(productivity proxies $\to$ input; soil-chemistry proxies $\to$ stabilisation).
Frame as ``consistent with'', not proof. Expect modest holdout gains against a low
predictability ceiling --- modest-but-robust-and-honest beats an overclaimed jump.
}
\fi
% =====================================================================
%  end optional strand
% =====================================================================

\section{Conclusions}
% Written LAST; author shapes. Closes on: success -> Bayesian enabling ->
% refine with historical data collection.

\subsection{The experiment succeeded}
Transient, calibrated initialization let a model ladder reproduce the observed
accumulation and its incipient saturation --- the dynamic an equilibrium start
forecloses. State the proof of concept plainly.

\subsection{The Bayesian approach made it possible}
% [Bayes]
Treating the historical start as a calibrated parameter, with uncertainty
propagated into state and dynamics, is what turned the idea into a working
implementation. Credit the machinery without overclaiming novelty.

\subsection{Refinement needs historical data}
% [History]
The approach can go much further once the pre-observation history is
reconstructed from data rather than left to calibration --- the open frontier,
and the natural successor study. Close here.

% =====================================================================
\part{Appendices}
% Modular methods inserts. Each \subfile below is a standalone `subfiles`
% document under appendices/ that also compiles on its own. Comment a line to
% drop that module from the combined build; the full sigma_input module is
% written, the rest are structural stubs to be filled.
% =====================================================================
\appendix

\subfile{appendices/appendix_sigma_input}          % FULL — physically-bounded input priors
\subfile{appendices/appendix_prior_homogenization} % stub
\subfile{appendices/appendix_transient_init}       % stub
\subfile{appendices/appendix_nonidentifiability}   % stub
\subfile{appendices/appendix_exact_integration}    % stub

% ---------------------------------------------------------------------
\section*{Notes for collapsing (remove before submission)}
\begin{itemize}
  \item \textbf{Parts} = OCAR acts; will dissolve into standard
        Intro/Methods/Results/Discussion at cleanup.
  \item Many heading levels are intentional scaffolding; merge aggressively
        once the story is fixed.
  \item Keep the protagonist (initialization problem) constant; do not let
        $\sigma_{\mathrm{input}}$ or non-identifiability take over the spine.
\end{itemize}

\end{document}
